% Einheit 3 – Sachaufgaben – Hauptnummern 40 bis 46
\einheitenkopf[e3][3 Sachaufgaben]{Einheit 3 von 4 · Sachaufgaben}
\zweigzeile{Hier lernst du, Zahlenrätsel und Flächenaufgaben mit quadratischen Gleichungen zu lösen · neu in diesem Jahr, je nach Buch bis Klasse 10 · keine P10-Aufgabe · baut auf: Wurzelziehen (Einheit 1), p-q-Formel (Einheit 2)}
\setcounter{aufgabe}{39}

\begin{aufgabe}{Ich erkenne, welche Lösung zur Frage passt.}
Die Gleichung ist schon gelöst. Kreuze an, welche Lösung als Antwort passt. Rechne nichts.
\begin{teile}
\teil Gesucht ist die Seitenlänge eines Quadrats. Lösungen: $9$ und $-9$ \hfill \kreuz{$9$} \kreuz{$-9$} \kreuz{beide}
\teil Gesucht ist eine Zahl. Lösungen: $4$ und $-4$ \hfill \kreuz{$4$} \kreuz{$-4$} \kreuz{beide}
\teil Gesucht ist die Breite eines Rechtecks. Lösungen: $6$ und $-10$ \hfill \kreuz{$6$} \kreuz{$-10$} \kreuz{beide}
\teil Gesucht ist die Zahl der Sitzreihen im Kino. Lösungen: $-15$ und $12$ \hfill \kreuz{$-15$} \kreuz{$12$} \kreuz{beide}
\teil Gesucht ist eine natürliche Zahl. Lösungen: $7$ und $-8$ \hfill \kreuz{$7$} \kreuz{$-8$} \kreuz{beide}
\end{teile}
\end{aufgabe}

\verfahren{Zahlenrätsel}
\begin{aufgabe}{Ich kann ein Zahlenrätsel mit dem Quadrat einer Zahl lösen.}
Löse die Gleichung und gib alle passenden Zahlen an.
\begin{teile}
\teil Das Quadrat einer Zahl ist 169. \quad $x^2 = 169$ \hfill Zahlen: \leerfeld
\teil Das Quadrat einer Zahl ist 400. \quad $x^2 = 400$ \hfill Zahlen: \leerfeld
\teil Verdoppelt man das Quadrat einer Zahl, erhält man 128. \quad $2x^2 = 128$ \hfill Zahlen: \leerfeld
\teil Das Quadrat einer Zahl, vermindert um 10, ist 90. \quad $x^2 - 10 = 90$ \hfill Zahlen: \leerfeld
\end{teile}
\end{aufgabe}

\begin{aufgabe}{Ich kann zu einem Zahlenrätsel selbst die Gleichung aufstellen.}
Stelle eine Gleichung auf, löse sie und gib die passenden Zahlen an.
\begin{teile}
\teil Addiert man zum Quadrat einer Zahl 15, erhält man 96. \\ Gleichung: \leerfeld \quad Zahlen: \leerfeld
\teil Das Dreifache des Quadrats einer Zahl ist 75. \\ Gleichung: \leerfeld \quad Zahlen: \leerfeld
\teil Das Produkt einer Zahl mit ihrem Nachfolger ist 56. \\ Gleichung: \leerfeld \quad Zahlen: \leerfeld
\teil Das Produkt zweier aufeinanderfolgender natürlicher Zahlen ist 132. \\ Gleichung: \leerfeld \quad Zahlen: \leerfeld
\end{teile}
\end{aufgabe}

\verfahren{Flächenaufgaben}
\begin{aufgabe}{Ich kann eine Rechteckaufgabe mit einer quadratischen Gleichung lösen.}
Ein rechteckiger Spielplatz ist $3\,\text{m}$ länger als breit. Seine Fläche beträgt $40\,\text{m}^2$.

\rechteck[seiten={x+3,x,x+3,x},punkte=]{4}{2.5}

\begin{teile}
\teil Stelle eine Gleichung für die Fläche auf.
\teil Löse die Gleichung.
\teil Welche Lösung passt zur Frage? Begründe.
\teil Gib Breite und Länge des Spielplatzes an und mache die Probe.
\end{teile}
\end{aufgabe}

\begin{aufgabe}{Ich kann eine Flächenaufgabe mit Klammern lösen.}
Stelle eine Gleichung auf, löse sie und antworte mit Einheit.
\begin{teile}
\teil Ein quadratisches Foto steckt in einem Rahmen, der ringsum $2\,\text{cm}$ breit ist. Foto und Rahmen zusammen haben eine Fläche von $196\,\text{cm}^2$. Wie lang ist eine Seite des Fotos?
\end{teile}

\quadratrand{x}{2 cm}

\begin{teile}
\teil Ein quadratischer Garten wird auf einer Seite um $4\,\text{m}$ verlängert und auf der anderen Seite um $2\,\text{m}$ verkürzt. Das neue Rechteck hat eine Fläche von $55\,\text{m}^2$. Wie lang war eine Seite des Gartens?
\end{teile}

\rechteck[seiten={x+4,x-2,x+4,x-2},punkte=]{4}{2}

\end{aufgabe}

\begin{aufgabe}{Ich finde den Fehler in einer Sachaufgabe.}
Finde den Fehler, benenne ihn und korrigiere die Rechnung.
\begin{teile}
\teil Ein Rechteck ist $5\,\text{cm}$ länger als breit und hat eine Fläche von $24\,\text{cm}^2$. Tim rechnet: \rechnung{x \cdot (x + 5) &= 24 \\ x^2 + 5x - 24 &= 0 \\ x_1 &= 3, \quad x_2 = -8} Tims Antwort: „Die Breite ist $3\,\text{cm}$ oder $-8\,\text{cm}$.“
\teil Ein Rechteck ist $2\,\text{cm}$ länger als breit und hat eine Fläche von $48\,\text{cm}^2$. Mia rechnet: \rechnung{2x + 2 \cdot (x + 2) &= 48 \\ 4x + 4 &= 48 \\ x &= 11}
\end{teile}
\end{aufgabe}

\begin{aufgabe}{Ich kann begründen, welche Lösungen im Sachzusammenhang gelten.}
Begründe.
\begin{teile}
\teil Warum gilt bei der Seitenlänge eines Rechtecks nur die positive Lösung?
\teil Warum sind beim Rätsel „Das Quadrat einer Zahl ist 225.“ beide Lösungen richtig?
\end{teile}
\end{aufgabe}
